\documentclass[11pt]{ctexart}
%=====================================================================
% theme-violet.tex —— 讲义精讲 · 紫花主题
% 用法：\documentclass[11pt]{ctexart}
%       \input{assets/theme-violet}
%       \begin{document} ... \end{document}
% 只用 ctexart 自带能力 + xcolor/enumitem/fancyhdr/pifont/amsmath，
% 不依赖 tikz/tcolorbox，保证 MiKTeX 可直接编译。
%=====================================================================

\usepackage[a4paper,left=2.0cm,right=2.0cm,top=2.1cm,bottom=2.3cm,%
            headsep=12pt,footskip=22pt]{geometry}
\usepackage{amsmath,amssymb}
\usepackage{needspace}
\usepackage{xcolor}
\usepackage{array}
\usepackage{enumitem}
\usepackage{pifont}
\usepackage{fancyhdr}
\usepackage[hidelinks]{hyperref}

%---------------------------------------------------------------- 配色
\definecolor{violet}{RGB}{107,78,143}      % 主色：紫
\definecolor{violetdeep}{RGB}{63,44,92}    % 深紫：标题/强调
\definecolor{petal}{RGB}{183,159,208}      % 花瓣：装饰线、花饰
\definecolor{bandviolet}{RGB}{247,243,251} % 浅紫底：档案/提示框
\definecolor{bandcream}{RGB}{252,249,242}  % 米白底：易错/注意框
\definecolor{rulepurple}{RGB}{217,207,230} % 分隔线
\definecolor{inkgray}{RGB}{92,96,110}      % 次要文字

%---------------------------------------------------------------- 字体
% 有 Noto CJK 用 Noto，没有就退回 ctex 默认（Windows 下为宋体/黑体）
\IfFontExistsTF{Noto Serif CJK SC}{\setCJKmainfont{Noto Serif CJK SC}}{}
\IfFontExistsTF{Noto Sans CJK SC}{\setCJKsansfont{Noto Sans CJK SC}}{}

%---------------------------------------------------------------- 版式
\setlist{nosep,leftmargin=1.9em}
\linespread{1.28}
\setlength{\parskip}{3.5pt}
\setlength{\parindent}{0pt}
\setlength{\tabcolsep}{5pt}
\renewcommand{\arraystretch}{1.28}

%---------------------------------------------------------------- 花饰
% \ding{95} = 五瓣小花，\ding{96} = 星形花；都来自 pifont，必定存在
\newcommand{\orn}{\textcolor{petal}{\ding{95}}}          % 行内小标
\newcommand{\ornstar}{\textcolor{petal}{\ding{96}}}      % 星形花
\newcommand{\flowerrule}{%                                % 居中花饰分隔线
  \par\vspace{7pt}%
  \noindent\makebox[\linewidth]{\color{petal}%
    \rule{0.20\linewidth}{0.6pt}\hspace{8pt}{\small\ding{95}}\hspace{4pt}%
    \raisebox{1.1pt}{\tiny\ding{95}}\hspace{4pt}{\small\ding{95}}\hspace{8pt}%
    \rule{0.20\linewidth}{0.6pt}}%
  \par\vspace{7pt}}

%---------------------------------------------------------------- 标题与节头
\newcommand{\doctitle}[3]{%   #1 主标题  #2 副标题  #3 元信息（可用 \\ 分行）
  \begin{center}
    {\color{petal}\small \orn\enspace #2\enspace \ornstar}\\[7pt]
    {\LARGE\bfseries\color{violetdeep} #1}\\[10pt]
    {\color{petal}\rule{0.55\linewidth}{1pt}}\\[4pt]
    {\small\color{inkgray} #3}
  \end{center}\vspace{2pt}}

\newcommand{\hdr}[2]{%   #1 章节号（如 §0）  #2 标题
  \par\vspace{12pt}%
  \noindent\colorbox{violetdeep}{\color{white}\bfseries\small\;#1\;}\hspace{7pt}%
  {\color{violetdeep}\large\bfseries #2}\par\nobreak
  \vspace{3pt}{\color{petal}\hrule height 1.2pt}\vspace{6pt}}

\newcommand{\subhdr}[1]{%
  \par\vspace{6pt}%
  \noindent{\color{violet}\bfseries #1}\par\vspace{1pt}}

%---------------------------------------------------------------- 提示框
\newcommand{\callout}[1]{% 浅紫底：讲义档案、读法说明
  \par\medskip\noindent
  \fcolorbox{rulepurple}{bandviolet}{\parbox{\dimexpr\linewidth-2\fboxsep-2\fboxrule}{\vspace{2pt}#1\vspace{2pt}}}%
  \par\medskip}

\newcommand{\warnbox}[1]{% 米白底：易错、边界条件
  \par\medskip\noindent
  \fcolorbox{rulepurple}{bandcream}{\parbox{\dimexpr\linewidth-2\fboxsep-2\fboxrule}{\vspace{2pt}#1\vspace{2pt}}}%
  \par\medskip}

\newcommand{\gapitem}[2]{% #1 缺口类型（紫） #2 内容
  \par\vspace{2pt}\noindent
  \hangindent=1.4em\hangafter=1
  {\color{petal}\small\ding{96}}\enspace{\bfseries\color{violetdeep}#1}\enspace #2\par}

% 题干块（\qstem）：每道例题／作业的「详解」开头必须先复述题干与出处。
% 起因：2026-10-25 第五章讲义把七道作业的题面集中放在一页、详解另起一页，
% 读者翻到「第 4 题」时看不到条件，等于一道没有题干的题目。
% 规则：详解块内自带题干，读者不翻页就能读懂题在问什么。
\newcommand{\qstem}[1]{%
  \par\medskip\noindent
  \fcolorbox{violet}{bandviolet}{%
    \parbox{\dimexpr\linewidth-2\fboxsep-2\fboxrule}{%
      \vspace{2pt}%
      {\small\bfseries\color{violet}题干}\hspace{0.8em}#1%
      \vspace{2pt}}}%
  \par\medskip}

% v2: 保护题干与开头，长题可跨页，续页补上下文。
\newcommand{\qbreak}{\needspace{6\baselineskip}}
% 独立练习：ID、来源、题干、考点、答案标签（通常 ans:ID）。
\newcommand{\practice}[5]{%
  \qbreak\subhdr{练习 #1}\qstem{（#2）#3}%
  \textbf{考点}：#4\par
  详解见附录第~\pageref{#5}~页。\par\medskip}
% 附录开头：ID、来源、条件摘要；自动产生 ans:ID 标签。
\newcommand{\answerstart}[3]{%
  \qbreak\subhdr{#1\quad 详解}\label{ans:#1}\qstem{（#2）#3}}
\newcommand{\answercontinue}[2]{%
  \needspace{5\baselineskip}\subhdr{#1\quad 续解}%
  {\small\color{inkgray}本页接续：#2}\par\medskip}
% 延迟提供默认值，让旧讲义可在导言区保留自己的 \newcommand{\step}。
\AtBeginDocument{\providecommand{\step}[1]{\needspace{3\baselineskip}\par\medskip\textbf{#1}\par}}

% 前提旁注：只放在「需要理解的那一步」后面，给陈述 + 本题代入，不重证。
\newcommand{\aside}[1]{%
  \par\vspace{1pt}\noindent
  \hangindent=1.6em\hangafter=1
  {\small\color{inkgray}%
    \hspace*{0.6em}{\color{petal}$\triangleright$}\hspace{2pt}#1}%
  \par\vspace{2pt}}

%---------------------------------------------------------------- 表格
% 本机 MiKTeX 的 colortbl 与 array 版本不匹配（\rowcolor 会报错），
% 故表格统一用「上粗紫线 + 表头 + 中细线」的三线式，不依赖 colortbl。
\newcommand{\tabtop}{\par\vspace{3pt}{\color{violet}\hrule height 1.2pt}\vspace{1pt}}
\newcommand{\tabmid}{\par\vspace{1pt}{\color{petal}\hrule height 0.5pt}\vspace{1pt}}
\newcommand{\tabbot}{\par\vspace{3pt}{\color{rulepurple}\hrule height 0.5pt}\vspace{8pt}}
\newcommand{\thead}[1]{{\bfseries\color{violetdeep}#1}}

%---------------------------------------------------------------- 页眉页脚
\pagestyle{fancy}
\fancyhf{}
\renewcommand{\headrulewidth}{0.4pt}
\renewcommand{\headrule}{\hbox to\headwidth{\color{rulepurple}\leaders\hrule height \headrulewidth\hfill}}
\fancyhead[L]{\small\color{inkgray}\songti \doctag}
\fancyhead[R]{\small\color{inkgray}\songti \coursetag}
\fancyfoot[C]{\small\color{inkgray}\orn\enspace{\color{violet}\thepage}\enspace\orn}

% 每篇讲义在导言区设这两个值
\newcommand{\doctag}{}
\newcommand{\coursetag}{}
\newcommand{\settags}[2]{\renewcommand{\doctag}{#1}\renewcommand{\coursetag}{#2}}

%---------------------------------------------------------------- 数学记号
\newcommand{\vect}[1]{\boldsymbol{#1}}
\newcommand{\dd}{\,\mathrm{d}}
\newcommand{\ddx}{\frac{\mathrm{d}}{\mathrm{d}x}}
\addtolength{\jot}{2pt}



% theme 的 \tabmid 在 tabular 内会把横线塞进首格，这里改成 \noalign 形式
\renewcommand{\tabmid}{\noalign{\vskip 2pt{\color{petal}\hrule height 0.5pt}\vskip 3pt}}
\settags{Polarization · 作业与速查}{光学（02）· 外教 Lecture 1}
\begin{document}

\doctitle{外教课 Polarization：作业解答与考前速查}
{2026-11-09（周一）课前准备}
{依据：\texttt{Polarization.pptx}（69 页，Lecture 1）与 \texttt{Homework Task.docx}\\
光学（02）｜整理日期 2026-11-09}

%=====================================================================
\hdr{§0}{今天做什么、明天交什么}

这门外教课只有三条硬规矩，弄错一条就白做：\textbf{手写}、\textbf{格式齐全}
（Given／公式／步骤／带单位的最终答案）、\textbf{英文签名}。其余都好办。

\small
\tabtop
\begin{tabular}{@{}p{3.4cm}p{5.6cm}p{6.2cm}@{}}
\thead{时间} & \thead{要做什么} & \thead{硬性要求} \\
\tabmid
今天 11-09 上课 & 外教讲 Polarization（Lecture 1） & 课前不用交作业（这是第一讲），带纸笔记要点 \\
\tabmid
今天 11-09 课末 10 分钟 & 小考：5 题，每题 4 个选项 & 允许使用任何资料；卷面签英文名；考完放讲台 \\
\tabmid
明天 11-10 课前 & 交本次布置的 3 道作业（共 30 分，每题 10 分） & 手写、步骤清楚、英文签名，开课前放到讲台 \\
\end{tabular}
\tabbot
\normalsize

\callout{本文件有三块，可以分开用：\textbf{§2 是考场速查卡}，直接打印带进考场（外教明确允许用任何资料）；
\textbf{§3 是三道作业的完整解答}，照抄成手写稿即可交；\textbf{§4 是十道模拟选择题}，用来验一验自己会不会。}

%=====================================================================
\hdr{§1}{三句话先把物理讲明白}

\textbf{第一句：光是横波。}电场 $\vect E$ 和磁场 $\vect B$ 都垂直于传播方向，也互相垂直。
真正跟物质打交道的是电场（它推得动电子），所以"偏振"说的永远是 $\vect E$ 的朝向。

\textbf{第二句：偏振＝$\vect E$ 在"垂直于光的那个平面"里怎么摆。}
固定指向一条线，叫线偏振；匀速转圈且长度不变，叫圆偏振；越转越长越短地描椭圆，叫椭圆偏振；
每一时刻方向都随机、各个方向机会均等，叫自然光（非偏振光）。

\textbf{第三句：产生偏振的办法只有一类逻辑——让某一个方向的 $\vect E$ 过不去。}
吸收（二向色性／偏振片）、反射与折射（Brewster 角）、散射（蓝天）、双折射（方解石、波片），
四种方法各走一条路，但都是在"挑方向"。

%=====================================================================
\hdr{§2}{考场速查卡（可打印带入）}

外教说了"任何资料都可以用"，所以这张卡就是你的主要武器。三张表加五个坑，五分钟能扫完。

\subhdr{表 1｜会算的六个式子，以及"用之前必须确认什么"}

\small
\tabtop
\begin{tabular}{@{}p{4.4cm}p{4.9cm}p{5.6cm}@{}}
\thead{情形} & \thead{式子} & \thead{用之前必须确认} \\
\tabmid
自然光过第一片理想偏振片 & $I_1=I_0/2$ & 入射必须是自然光；偏振片无吸收损失 \\
\tabmid
线偏振光过检偏器（Malus 定律） & $I=I_{\text{前}}\cos^2\theta$ & 入射已经是\emph{线偏振}；$\theta$ 是入射 $\vect E$ 与透振方向的夹角 \\
\tabmid
偏振度 & $P=\dfrac{I_{\max}-I_{\min}}{I_{\max}+I_{\min}}$ & 转动检偏器读出两个极值；$P=1$ 是全偏振，$P=0$ 是自然光 \\
\tabmid
Brewster 角 & $\tan\theta_B=n_{21}=\dfrac{n_2}{n_1}$ & 界面是电介质（水、玻璃）；金属不成立 \\
\tabmid
Brewster 角的配套关系 & $\theta_B+\theta_2=90^\circ$ & 只在 Brewster 角成立，别处不成立 \\
\tabmid
s 分量的反射率 & $R_s=\sin^2(\theta_B-\theta_2)$ & 只在 Brewster 角用；此时 $R_p=0$ \\
\tabmid
Fresnel 振幅系数（s 分量） & $r_s=\dfrac{n_1\cos\theta_1-n_2\cos\theta_2}{n_1\cos\theta_1+n_2\cos\theta_2}$ & 光强反射率是 $R=|r|^2$，别把 $r$ 当 $R$ \\
\tabmid
Fresnel 振幅系数（p 分量） & $r_p=\dfrac{n_2\cos\theta_1-n_1\cos\theta_2}{n_2\cos\theta_1+n_1\cos\theta_2}$ & 分子里的 $n_1$、$n_2$ 与 s 分量正好互换 \\
\tabmid
波片的相位差 & $\Delta\varphi=\dfrac{2\pi}{\lambda}(n_o-n_e)d$ & 厚度 $d$ 是光在晶体里走过的距离 \\
\tabmid
四分之一波片 & $d=\dfrac{\lambda}{4\,|n_o-n_e|}$ & 正晶体 $n_e>n_o$，负晶体反过来，所以取绝对值 \\
\tabmid
半波片 & $d=\dfrac{\lambda}{2\,|n_o-n_e|}$ & 正好是四分之一波片厚度的两倍 \\
\end{tabular}
\tabbot
\normalsize

\needspace{7cm}
\subhdr{表 2｜四种产生偏振的方法（选择题最爱考"这是哪一种"）}

\small
\tabtop
\begin{tabular}{@{}p{3.4cm}p{7.0cm}p{4.4cm}@{}}
\thead{方法} & \thead{一句话原理} & \thead{代表器件／现象} \\
\tabmid
选择性吸收（二向色性） & 晶体对两个互相垂直的振动方向吸收本领不同，一个方向被吃掉，另一个留下来 & Polaroid 偏振片（PVA 拉伸＋碘）、电气石、碘硫酸奎宁 \\
\tabmid
反射与折射 & 在 Brewster 角上 p 分量完全不能反射，反射光只剩 s 分量，是 100\% 线偏振 & 水面与玻璃的反光、偏振太阳镜、摄影偏振镜 \\
\tabmid
散射 & 空气分子里的电子被 $\vect E$ 推着振荡，像小天线那样再辐射；沿 90$^\circ$ 看只剩一个分量 & 蓝天（偏振度最高 85--95\%）、蜜蜂与蚂蚁的导航 \\
\tabmid
双折射 & 各向异性晶体把一束光劈成 o 光与 e 光，两者偏振方向互相垂直 & 方解石（冰洲石）、Nicol 棱镜、Glan--Taylor 棱镜、$\lambda/4$ 与 $\lambda/2$ 波片 \\
\end{tabular}
\tabbot
\normalsize

\needspace{6cm}
\subhdr{表 3｜相位差决定出射是什么偏振}

\small
\tabtop
\begin{tabular}{@{}p{3.2cm}p{11.2cm}@{}}
\thead{相位差 $\Delta\varphi$} & \thead{出射光的偏振态} \\
\tabmid
$0$（或 $2\pi$ 的整数倍） & 线偏振，方向不变。这就是全波片 \\
\tabmid
$\pi/2$（$90^\circ$） & 入射线偏振光若与快轴成 $45^\circ$，出射是圆偏振。这就是 $\lambda/4$ 波片 \\
\tabmid
$\pi$（$180^\circ$） & 出射仍是线偏振，但方向相对主轴转过了 $2\alpha$（$\alpha$ 是入射方向与主轴的夹角）。这就是 $\lambda/2$ 波片 \\
\tabmid
其他值 & 椭圆偏振 \\
\end{tabular}
\tabbot
\normalsize

\subhdr{五个最容易丢分的坑}

\warnbox{
\textbf{坑 1：拿 Malus 定律直接乘自然光。}自然光 $1200\ \mathrm{W/m^2}$ 过第一片得到的是 $600$，
不是 $1200\cos^2 30^\circ=900$。Malus 定律的前提是"入射已经是线偏振"。

\textbf{坑 2：把角度相加。}三片偏振片时，$30^\circ$ 与 $45^\circ$ 是相邻两片之间的夹角，不是 $75^\circ$。
光过第二片之后，$\vect E$ 已经沿线二方向，所以第三片只跟第二片差 $45^\circ$。

\textbf{坑 3：Brewster 用 $\sin$ 不用 $\tan$。}最省事是 $\tan\theta_B=n_2/n_1$，
再用 $\theta_2=90^\circ-\theta_B$ 一步到位，不必先套 Snell 求折射角。

\textbf{坑 4：算反射光强时忘了 s 分量只占自然光的一半。}自然光里 s 与 p 各占 $I_0/2$，
反射光强是 $(I_0/2)R_s$，不是 $I_0R_s$。

\textbf{坑 5：波片厚度里的折射率差忘了绝对值。}石英是正晶体，$n_e>n_o$，
$n_o-n_e=-0.009$，而厚度公式里必须用 $|n_o-n_e|=0.009$。
}

%=====================================================================
\hdr{§3}{作业三道题：可直接照抄的完整解答}

下面每一题的英文块就是\textbf{照抄内容}（手写时把 $\Rightarrow$ 写成箭头、分数写成横线即可）；
英文块后面的"中文讲评"是给你自己看的，不要抄上去。

%---------------------------------------------------------------------
\qbreak
\subhdr{Problem 1｜Malus 定律与自然光}

\qstem{（\texttt{Homework Task.docx} Problem 1）Natural light of intensity $I_0=1200\ \mathrm{W/m^2}$
passes through three polaroids in series. The angle between the transmission planes of the first
and second polaroids is $30^\circ$, and the angle between the second and third is $45^\circ$.
Find (a) the intensity after each polaroid; (b) the total fraction of the original intensity
transmitted.}

\textbf{Given}\quad
$I_0=1200\ \mathrm{W/m^2}$ (unpolarized natural light);\;
$\theta_{12}=30^\circ$ (between polaroids 1 and 2);\;
$\theta_{23}=45^\circ$ (between polaroids 2 and 3).

\textbf{Formulas}\quad
Natural light through the first polaroid: $I_1=I_0/2$.\;
Malus's law for every later polaroid: $I=I_{\text{before}}\cos^2\theta$.

\textbf{Step 1. After the first polaroid.}\ Natural light has no preferred direction, so the first polaroid passes exactly one half:
\[
I_1=\frac{I_0}{2}=\frac{1200}{2}=600\ \mathrm{W/m^2},
\]
and the transmitted light is linearly polarized along the axis of polaroid 1.

\textbf{Step 2. After the second polaroid.}\ The incident light is now linearly polarized, so Malus's law applies with
$\theta_{12}=30^\circ$:
\[
I_2=I_1\cos^2 30^\circ=600\times\left(\frac{\sqrt3}{2}\right)^2=600\times\frac34
=450\ \mathrm{W/m^2}.
\]

\textbf{Step 3. After the third polaroid.}\ The light leaving polaroid 2 oscillates along the axis of polaroid 2, so the relevant angle is
$\theta_{23}=45^\circ$:
\[
I_3=I_2\cos^2 45^\circ=450\times\frac12=225\ \mathrm{W/m^2}.
\]

\textbf{Answer (a)}\quad
$I_1=600\ \mathrm{W/m^2}$,\quad $I_2=450\ \mathrm{W/m^2}$,\quad $I_3=225\ \mathrm{W/m^2}$.

\textbf{Answer (b)}\quad
\[
\frac{I_3}{I_0}=\frac{225}{1200}=0.1875=18.75\%,
\qquad\text{equivalently}\qquad
\frac{I_3}{I_0}=\frac12\cdot\frac34\cdot\frac12=\frac{3}{16}.
\]

\textbf{中文讲评。}这道题只有一个转折点：第一片是"减半"，后面两片才是"乘 $\cos^2$"。
物理上为什么减半——自然光可以看成两个互相垂直、强度相等、彼此不相干的线偏振分量之和，
偏振片只放过其中一个，所以剩一半。第二片之后光的振动方向已经跟着第二片走，
所以第三片只需要跟第二片比，题目给的 $30^\circ$ 到这一步就不再用得上了。

\textbf{自检。}最终比例 $0.1875$ 小于 $1/2$，符合"过第一片就损失一半"的直觉；
若你算出 $900\ \mathrm{W/m^2}$ 这种数，就是踩了坑 1。

%---------------------------------------------------------------------
\qbreak
\subhdr{Problem 2｜Brewster 角与 Fresnel 反射}

\qstem{（\texttt{Homework Task.docx} Problem 2）A beam of natural light is incident from air onto a
glass plate with refractive index $n=1.6$. Find (a) the Brewster angle; (b) the angle of refraction
at the Brewster angle; (c) the intensity reflection coefficient for s-polarized light at the Brewster
angle, using $R_s=\sin^2(\theta_B-\theta_2)$; (d) the intensity of the reflected light if the
incident natural light has intensity $I_0=500\ \mathrm{W/m^2}$.}

\textbf{Given}\quad
$n_1=1.00$ (air);\; $n_2=1.6$ (glass);\; $n_{21}=n_2/n_1=1.6$;\;
$I_0=500\ \mathrm{W/m^2}$ (natural light).

\textbf{Formulas}\quad
Brewster's law: $\tan\theta_B=n_{21}$.\;
At the Brewster angle the reflected and refracted rays are perpendicular, so
$\theta_B+\theta_2=90^\circ$.\;
$R_s=\sin^2(\theta_B-\theta_2)$;\; intensity reflectance $R=|r|^2$;\; $R_p=0$ at $\theta_B$.

\textbf{Step 1. Brewster angle.}\par\noindent\[
\tan\theta_B=n_{21}=1.6
\quad\Longrightarrow\quad
\theta_B=\arctan(1.6)=57.99^\circ\approx 58.0^\circ.
\]

\textbf{Step 2. Angle of refraction.}\ At the Brewster angle the reflected and refracted rays are perpendicular, hence
\[
\theta_2=90^\circ-\theta_B=90^\circ-57.99^\circ=32.01^\circ\approx 32.0^\circ .
\]
(Check with Snell's law: $1.00\sin 57.99^\circ=0.8480$ and $1.6\sin 32.01^\circ=0.8480$ — they agree.)

\textbf{Step 3. Reflection coefficient for s-polarized light.}\par\noindent\[
\theta_B-\theta_2=57.99^\circ-32.01^\circ=25.99^\circ,
\]
\[
R_s=\sin^2(25.99^\circ)=(0.4382)^2=0.192\approx 19.2\% .
\]
(Cross-check with the Fresnel formula:
$r_s=\dfrac{n_1\cos\theta_B-n_2\cos\theta_2}{n_1\cos\theta_B+n_2\cos\theta_2}
=\dfrac{0.5300-1.3568}{0.5300+1.3568}=-0.4382$, so $R_s=r_s^2=0.192$.)

\textbf{Step 4. Intensity of the reflected light.}\ Natural light contains equal s and p components, each carrying $I_0/2$. At the Brewster angle
$R_p=0$, so only the s component is reflected:
\[
I_{\text{refl}}=\frac{I_0}{2}\,R_s=\frac{500}{2}\times 0.192=48.0\ \mathrm{W/m^2}.
\]

\textbf{Answer}\quad
(a) $\theta_B\approx 58.0^\circ$;\quad
(b) $\theta_2\approx 32.0^\circ$;\quad
(c) $R_s\approx 0.192$ (19.2\%);\quad
(d) $I_{\text{refl}}\approx 48\ \mathrm{W/m^2}$, and this reflected beam is 100\% s-polarized
(oscillations perpendicular to the plane of incidence).

\textbf{中文讲评。}Brewster 角之所以特别，是因为此时 p 分量的振动方向正好沿着折射光的方向，
而横波不可能沿着自己的传播方向振动，所以 p 分量根本反射不出去，反射光里只剩 s 分量。
这就是"反射光是 100\% 线偏振"的物理原因，也是 (d) 里只乘 $R_s$、并且要先取 $I_0/2$ 的道理。

\textbf{自检。}$R_s=0.192$ 看着偏大，但 $\sin 26^\circ\approx 0.44$，平方就是 $0.19$，没有算错；
$R_s$ 随 $n$ 增大而增大：$n=1.33$ 时约 $0.077$、$n=1.5$ 时约 $0.15$、$n=2$ 时约 $0.36$，本问的 $0.19$ 正落在中间，数量级合理。

%---------------------------------------------------------------------
\qbreak
\subhdr{Problem 3｜四分之一波片}

\qstem{（\texttt{Homework Task.docx} Problem 3）A quarter-wave plate is made from quartz with
$n_o=1.544$ and $n_e=1.553$. The wavelength of the incident light is $\lambda=632.8\ \mathrm{nm}$
(He--Ne laser). Find (a) the minimum thickness required to convert linearly polarized light into
circularly polarized light; (b) if the thickness is doubled, what type of plate does it become
($\lambda/4$, $\lambda/2$, or full-wave)? Explain.}

\textbf{Given}\quad
$n_o=1.544$;\; $n_e=1.553$;\; $\lambda=632.8\ \mathrm{nm}=632.8\times10^{-9}\ \mathrm{m}$.

\textbf{Formulas}\quad
Phase difference introduced by a plate of thickness $d$:
$\Delta\varphi=\dfrac{2\pi}{\lambda}(n_o-n_e)d$.\;
For a quarter-wave plate $\Delta\varphi=\pi/2$, which gives
$d=\dfrac{\lambda}{4\,|n_o-n_e|}$.

\textbf{Step 1. The condition for circular polarization.}\ Linear input becomes circular output when the plate delays one component by a quarter of a
wavelength, i.e.\ $\Delta\varphi=\pi/2$. In units of wavelength this means
$(n_o-n_e)d=\lambda/4$, so there is exactly one unknown, $d$.

\textbf{Step 2. Solve for the minimum thickness.}\par\noindent\[
d=\frac{\lambda}{4\,|n_o-n_e|}=\frac{632.8\ \mathrm{nm}}{4\times|1.544-1.553|}
=\frac{632.8\ \mathrm{nm}}{4\times 0.009}=\frac{632.8}{0.036}\ \mathrm{nm}
=1.7578\times10^{4}\ \mathrm{nm}.
\]
\[
d=1.76\times10^{-5}\ \mathrm{m}\approx 17.6\ \mu\mathrm{m}.
\]

\textbf{Answer (a)}\quad $d_{\min}\approx 1.76\times10^{-5}\ \mathrm{m}\approx 17.6\ \mu\mathrm{m}$
(about $17\,578\ \mathrm{nm}$).

\textbf{Step 3. Double the thickness.}\ If $d'=2d$, the retardation doubles as well:
\[
(n_o-n_e)d'=2\times\frac{\lambda}{4}=\frac{\lambda}{2}
\quad\Longrightarrow\quad
\Delta\varphi'=2\times\frac{\pi}{2}=\pi .
\]
A phase difference of $\pi$ means the plate is a \textbf{half-wave ($\lambda/2$) plate}.
It does \emph{not} make the light circularly polarized: light that enters linearly polarized
leaves linearly polarized, with its plane of polarization rotated by $2\alpha$, where $\alpha$ is
the angle between the incident vibration direction and the fast axis of the plate.

\textbf{Answer (b)}\quad Doubling the thickness turns the $\lambda/4$ plate into a
$\lambda/2$ (half-wave) plate, because the retardation becomes $\lambda/2$ and the phase difference
becomes $\pi$.

\textbf{中文讲评。}波片不改变光的强度，它只让两个互相垂直的分量"错开一段路程"。
这段光程差 $\Delta L=(n_o-n_e)d$ 除以波长，就是差了几个周期；
$\Delta L=\lambda/4$ 时两分量差 $90^\circ$，圆偏振；$\Delta L=\lambda/2$ 时差 $180^\circ$，
线偏振被"翻面"，方向转过 $2\alpha$。所以厚度翻倍，正好从四分之一波片变成半波片。

\textbf{自检。}$17.6\ \mu\mathrm{m}$ 只有头发丝直径的五分之一左右，很薄但做得出来——
这正是石英波片真实的数量级；若你算出毫米级，肯定是把 $10^{-9}$ 与 $\mu\mathrm{m}$ 混了。

%=====================================================================
\clearpage
\hdr{§4}{模拟小测：十道选择题}

外教的卷子是 5 题、每题 4 选 1、允许翻资料，所以真正的难点是\textbf{读得快、找得到}。
下面十题按课件内容出，限时 12 分钟做完，答案在末尾。

\small
\begin{enumerate}[label=\textbf{Q\arabic*.},leftmargin=2.5em,itemsep=2pt,topsep=2pt]
\item Unpolarized light of intensity $I_0$ passes through a single ideal polarizer.
The transmitted intensity is
\begin{enumerate}[label=(\Alph*),leftmargin=2.4em,itemsep=0pt,topsep=1pt,parsep=0pt]
\item $I_0$ \qquad (B) $I_0/4$ \qquad (C) $I_0/2$ \qquad (D) $0$
\end{enumerate}

\item Linearly polarized light of intensity $I_0$ passes through an analyzer whose transmission
axis makes $60^\circ$ with the plane of polarization. The transmitted intensity is
\begin{enumerate}[label=(\Alph*),leftmargin=2.4em,itemsep=0pt,topsep=1pt,parsep=0pt]
\item $I_0/4$ \qquad (B) $I_0/2$ \qquad (C) $3I_0/4$ \qquad (D) $0$
\end{enumerate}

\item The Brewster angle for light going from air ($n=1$) into water ($n=1.33$) is closest to
\begin{enumerate}[label=(\Alph*),leftmargin=2.4em,itemsep=0pt,topsep=1pt,parsep=0pt]
\item $41.2^\circ$ \qquad (B) $48.8^\circ$ \qquad (C) $53.1^\circ$ \qquad (D) $56.3^\circ$
\end{enumerate}

\item At the Brewster angle, the reflected ray is
\begin{enumerate}[label=(\Alph*),leftmargin=2.4em,itemsep=0pt,topsep=1pt,parsep=0pt]
\item unpolarized
\item partially polarized
\item completely linearly polarized, with the electric field perpendicular to the plane of incidence
\item completely linearly polarized, with the electric field inside the plane of incidence
\end{enumerate}

\item At the Brewster angle, the sum $\theta_B+\theta_2$ of the angles of incidence and refraction equals
\begin{enumerate}[label=(\Alph*),leftmargin=2.4em,itemsep=0pt,topsep=1pt,parsep=0pt]
\item $90^\circ$ \qquad (B) $180^\circ$ \qquad (C) $\theta_B$ \qquad (D) it depends on the refractive index
\end{enumerate}

\item A quarter-wave plate converts linearly polarized light into circularly polarized light only when
\begin{enumerate}[label=(\Alph*),leftmargin=2.4em,itemsep=0pt,topsep=1pt,parsep=0pt]
\item the incident light is natural light
\item the incident vibration direction makes $45^\circ$ with the fast axis
\item the incident vibration direction is parallel to the fast axis
\item the plate is twice as thick as a half-wave plate
\end{enumerate}

\item A quarter-wave plate is cut from quartz ($n_o=1.544$, $n_e=1.553$) for $\lambda=600\ \mathrm{nm}$.
Its minimum thickness is closest to
\begin{enumerate}[label=(\Alph*),leftmargin=2.4em,itemsep=0pt,topsep=1pt,parsep=0pt]
\item $8.3\ \mu\mathrm{m}$ \qquad (B) $16.7\ \mu\mathrm{m}$ \qquad (C) $33.3\ \mu\mathrm{m}$ \qquad (D) $66.7\ \mu\mathrm{m}$
\end{enumerate}

\item Skylight is most strongly polarized when you look at the sky
\begin{enumerate}[label=(\Alph*),leftmargin=2.4em,itemsep=0pt,topsep=1pt,parsep=0pt]
\item straight towards the Sun
\item at $45^\circ$ from the Sun
\item at $90^\circ$ from the Sun
\item directly away from the Sun
\end{enumerate}

\item Polarizing sunglasses have a vertical transmission axis because glare reflected from
horizontal surfaces is
\begin{enumerate}[label=(\Alph*),leftmargin=2.4em,itemsep=0pt,topsep=1pt,parsep=0pt]
\item vertically polarized
\item horizontally polarized
\item circularly polarized
\item unpolarized
\end{enumerate}

\item When a linearly polarized beam passes through a half-wave plate, the emerging light is
\begin{enumerate}[label=(\Alph*),leftmargin=2.4em,itemsep=0pt,topsep=1pt,parsep=0pt]
\item circularly polarized
\item linearly polarized, with the plane rotated by $2\alpha$ ($\alpha$ = angle to the fast axis)
\item unpolarized
\item elliptically polarized
\end{enumerate}
\end{enumerate}
\normalsize

\needspace{8cm}
\subhdr{答案与一句解析}

\small
\tabtop
\begin{tabular}{@{}p{1.2cm}p{1.6cm}p{11.4cm}@{}}
\thead{题} & \thead{答案} & \thead{为什么} \\
\tabmid
Q1 & C & 自然光各方向均等，只放过一个分量，剩一半。课件 p60 第一行 \\
\tabmid
Q2 & A & $\cos^2 60^\circ=0.25$，Malus 定律。课件 p18--p19 \\
\tabmid
Q3 & C & $\tan\theta_B=1.33\Rightarrow\theta_B=53.1^\circ$。课件 p25 \\
\tabmid
Q4 & C & 反射光只剩 s 分量，s 垂直于入射面。课件 p31 FAQ 第一条 \\
\tabmid
Q5 & A & Brewster 角上反射光与折射光互相垂直。课件 p31、p53 \\
\tabmid
Q6 & B & 只有与快轴成 $45^\circ$ 时两分量振幅相等，才能合成圆偏振。课件 p50 \\
\tabmid
Q7 & B & $600/(4\times0.009)=16\,667\ \mathrm{nm}\approx16.7\ \mu\mathrm{m}$。课件 p58--p59 \\
\tabmid
Q8 & C & 90$^\circ$ 散射只剩一个分量，偏振度最高（85--95\%）。课件 p34 \\
\tabmid
Q9 & B & 水平面反射的眩光是水平偏振，竖直透振方向把它挡掉。课件 p31 FAQ 第二条 \\
\tabmid
Q10 & B & 半波片差 $180^\circ$，线偏振进、线偏振出，方向转 $2\alpha$。课件 p50 \\
\end{tabular}
\tabbot
\normalsize

%=====================================================================
\hdr{§5}{这份材料的口径说明}

\textbf{来源。}三张速查表按课件 p60「Summary of Key Formulas」逐条对齐；
s／p 的定义取课件 p27（s 来自德语 \emph{senkrecht}，p 来自 \emph{parallel}）；
四种产生方法的结构取课件 p12／p23／p32／p40 的四张"Method $n$ of 4"总览页。

\textbf{作业与课件例题的关系。}这三道作业不是新题型，而是课件 p52、p53、p56、p58 四道例题换了数字：
Problem 1 对应 p55「三片偏振片」；Problem 2 对应 p53 与 p56；Problem 3 对应 p58，只把 $\lambda$ 由
$589\ \mathrm{nm}$ 换成 $632.8\ \mathrm{nm}$。所以判分口径应当与课件示例一致：写 Given、写公式、
代数字、给带单位的答案。

\textbf{模拟题的可靠性。}§4 的十道题是按课件正文推出来的，\emph{不是}往年真题；
外教实际出题范围若与课件有出入，以课上讲的为准。

\textbf{没有写进能力账本。}这一份属于课前准备与作业辅助，没有产生你本人的独立作答记录，
所以它不会让任何考点变成"已掌握"。真要留下证据，可以等课后再做一份换过数字的题，
闭卷、不翻这份材料。

\flowerrule
{\small\color{inkgray}
一句话收尾：会算的三件事——自然光减半、线偏振乘 $\cos^2\theta$、波片看光程差是 $\lambda/4$ 还是 $\lambda/2$；
会背的三件事——$\tan\theta_B=n_{21}$、Brewster 角上反射光纯 s、90$^\circ$ 散射偏振最强。}

\end{document}

